\documentclass[10pt,a4paper]{amsart}
\usepackage[margin=19mm]{geometry}
\usepackage{amsmath,amssymb,mathtools}
\usepackage[scr=rsfs,cal=euler]{mathalfa}
\usepackage{xfrac}
\usepackage[unicode,hidelinks,bookmarks=false]{hyperref}
\newcommand{\ie}{i.e.\ }
\newcommand{\cf}{cf.\ }
\newcommand{\resp}{respectively\ }
\newcommand{\Subsection}{\subsection}
\providecommand{\nobreakdash}{\nobreak\hbox{-}}
\setlength{\emergencystretch}{2em}
\multlinegap0pt
\usepackage{bm}
\usepackage{bbm}
\usepackage{dsfont}
\usepackage{xspace}
\usepackage{cancel}
\usepackage{mathtools}
\usepackage{amsthm}
\makeatletter
\@ifundefined{thm}{\newtheorem{thm}{Thm}}{}
\@ifundefined{lem}{\newtheorem{lem}[thm]{Lemma}}{}
\makeatother
\title[A Transverse Averaging Operator and Cohomology of Quotients ]{A Transverse Averaging Operator and Cohomology of Quotients by Homogeneous Non-closed Subgroups}
\author{Yi Lin}
\date{Source: arXiv:2604.17619v2 (CC BY 4.0)}
\begin{document}
\maketitle

\noindent\emph{Source and licence.} A Transverse Averaging Operator and Cohomology of Quotients by Homogeneous Non-closed Subgroups. Version arXiv:2604.17619v2 (CC BY 4.0), distributed under the Creative Commons Attribution 4.0 International licence, as recorded in the arXiv licence metadata for this exact version and shown on the abstract page https://arxiv.org/abs/2604.17619v2. Full rights notice: licenses/arxiv-2604.17619-CC-BY-4.0.txt.\par\smallskip

\noindent\textit{Editorial scope.} Counted unit: the Lem whose source label is \texttt{coefficient-complex} in the article (source0.tex lines 1136--1190, source label \texttt{coefficient-complex}), delivered together with the article's own complete original proof of it (source0.tex lines 1192--1244, one \texttt{proof} environment of 53 lines). No mathematics is written, repaired or re-derived: statement and proof are the article's own text line for line. Required context retained uncounted: none. Every definition, display and label of those lines is retained unchanged. Omitted from this package and not redistributed: 1--1135, 1191--1191, 1245--2044. Nothing inside a retained excerpt is omitted or altered: every retained body is delivered exactly as the article's lines are written, so no omission receipt of any kind accompanies this package. Citations: the retained excerpts carry no citation command at all, so no bibliography entry of the article is transcribed and no reference is redistributed. Reference adaptation: none was needed and none was made; every reference inside the delivered unit resolves inside it, so no unresolved reference or citation is produced. Printed slips kept literal: no printed word, symbol, hypothesis or number of the counted statement or of its proof was altered. Layout and class: the article's own class is \texttt{amsart} with packages including geometry, srcltx, eucal, xy, mathrsfs, xypic, amsfonts, amsmath, amsthm, bm, amssymb, tikz-cd, latexsym, tabularx; the wrapper is the lane's pinned \texttt{amsart} editorial wrapper at 10pt on A4, which restates the theorem-like environments the retained text needs and adds the article's own macro definitions that the retained text needs (none), with extra wrapper packages (bm, bbm, dsfont, xspace, cancel, mathtools). The delivered numbering is the unit's own. Assets: no retained span contains a figure environment, a graphics-inclusion command, a PDF-page inclusion, a table or a TikZ picture; no image file is copied into this package, the package declares no asset dependency and no image file is redistributed with it. Licence: version arXiv:2604.17619v2 of this article is distributed under the Creative Commons Attribution 4.0 International licence, as recorded in the arXiv licence metadata for this exact version and shown on the abstract page https://arxiv.org/abs/2604.17619v2. A full rights notice is delivered at licenses/arxiv-2604.17619-CC-BY-4.0.txt. Characters: every retained excerpt is pure ASCII, so the delivered engine, XeLaTeX with its default Latin Modern text font, reports no missing character.
\begin{lem}\label{coefficient-complex}
The following statements hold.
\begin{itemize}
\item[(a)] For every $k$, the assignment
\begin{equation}\label{coefficient-map}
F_\beta(p(x))(\xi_1,\dots,\xi_k)
:=
\beta_x(\overline X_{\xi_1},\dots,\overline X_{\xi_k})
\end{equation}
defines an isomorphism
\begin{equation}\label{coefficient-isomorphism}
\mathcal I_k:
\Omega^k(M,\mathcal F)
\xrightarrow{\;\cong\;}
C^\infty(B,\wedge^k\mathfrak g^*),
\qquad
\beta\longmapsto F_\beta.
\end{equation}

\item[(b)] Under \eqref{coefficient-isomorphism}, the exterior differential
corresponds to the operator
$
D:C^\infty(B,\wedge^k\mathfrak g^*)
\longrightarrow
C^\infty(B,\wedge^{k+1}\mathfrak g^*)$
given by
\begin{align}
(DF)(\xi_0,\dots,\xi_k)
={}&
\sum_{i=0}^k(-1)^i(\xi_i)_B
\bigl(
F(\xi_0,\dots,\widehat{\xi_i},\dots,\xi_k)
\bigr)
\nonumber\\
&+
\sum_{0\leq i<j\leq k}(-1)^{i+j}
F\bigl(
[\xi_i,\xi_j],
\xi_0,\dots,\widehat{\xi_i},\dots,
\widehat{\xi_j},\dots,\xi_k
\bigr).
\label{coefficient-differential}
\end{align}

\item[(c)] For every $\xi\in\mathfrak g$ and
$\beta\in\Omega^k(M,\mathcal F)$,
\begin{equation}\label{coefficient-lie-derivative}
F_{\mathcal L_{\rho(\xi)}\beta}
=
\xi_BF_\beta+\operatorname{ad}^*(\xi)F_\beta,
\end{equation}
where $\xi_BF_\beta$ denotes componentwise differentiation of the
$\wedge^k\mathfrak g^*$-valued function $F_\beta$ along $\xi_B$.
\end{itemize}
\end{lem}

\begin{proof}
For fixed $\xi_1,\dots,\xi_k$, the right-hand side of
\eqref{coefficient-map} is a basic function on $M$ and is therefore
constant on each leaf closure. Thus $F_\beta$ is well defined.

Conversely, if
$F\in C^\infty(B,\wedge^k\mathfrak g^*)$, define
\[
(\mathcal I_k^{-1}F)_x(v_1,\dots,v_k)
=
F(p(x))
\bigl(
\alpha_x(v_1),\dots,\alpha_x(v_k)
\bigr).
\]
This is a basic form, and the two constructions are inverse to each other.
This proves~(a).

Part~(b) follows from the formula for the exterior differential and the
identity
\[
[\overline X_\xi,\overline X_\eta]
=
\overline X_{[\xi,\eta]}.
\]
Indeed, evaluating $d\beta$ on
$\overline X_{\xi_0},\dots,\overline X_{\xi_k}$ gives precisely
\eqref{coefficient-differential}.

For~(c), fix $\xi,\xi_1,\dots,\xi_k\in\mathfrak g$, and choose foliate
representatives $X$ and $Y_i$ of $\overline X_\xi$ and
$\overline X_{\xi_i}$, respectively. The function
$\beta(Y_1,\dots,Y_k)$ is basic and descends to
$F_\beta(\xi_1,\dots,\xi_k)$ on $B$. Moreover, $X$ projects to $\xi_B$,
while the transverse class of $[X,Y_i]$ is
$\overline X_{[\xi,\xi_i]}$. Using the Leibniz rule for the Lie derivative,
we obtain
\begin{align*}
F_{\mathcal L_{\rho(\xi)}\beta}(\xi_1,\dots,\xi_k)
&=
X\bigl(\beta(Y_1,\dots,Y_k)\bigr)
-\sum_{i=1}^k
\beta(Y_1,\dots,[X,Y_i],\dots,Y_k)\\
&=
\xi_B\bigl(F_\beta(\xi_1,\dots,\xi_k)\bigr)
-\sum_{i=1}^k
F_\beta(\xi_1,\dots,[\xi,\xi_i],\dots,\xi_k)\\
&=
\bigl(\xi_BF_\beta+\operatorname{ad}^*(\xi)F_\beta\bigr)
(\xi_1,\dots,\xi_k).
\end{align*}
This proves~(c).
\end{proof}

\end{document}
